\chapter{\MakeUppercase{Conclusão}}\label{sec:conclusao}

Este trabalho construiu e avaliou o AgroVisible, uma ferramenta que acopla uma rede LSTM alimentada por variáveis exógenas a um processo de decisão de Markov resolvido por programação dinâmica para prescrever quando vender o café arábica colhido. A avaliação seguiu um protocolo de estudo empírico, com questões de pesquisa definidas a priori, partição cronológica estrita, \textit{backtest} em safras cegas e análise de sensibilidade.

\section{Síntese das respostas}

\paragraph{RQ1.} A rede, na configuração final de 32 e 16 unidades com \textit{dropout} 0,4, escolhida por busca em validação e verificada com cinco sementes, não supera o passeio aleatório no horizonte de 21 dias e perde para a curva de futuros no de 63 dias na validação. No horizonte de 126 dias ela vence todas as referências nos dois conjuntos, e no teste cego vence as duas referências em 63 e 126 dias, com acerto direcional de 76 e 77\,\%. O valor preditivo está na direção de médio prazo, não na magnitude nem no curto prazo. A análise exploratória, a importância por permutação e a ablação convergem: o sinal vem da oferta, da base físico-futuro, do calendário e da bienalidade, cresce com o horizonte, e as variáveis de curto prazo são ruído. Cinco dos seis grupos contribuem no horizonte de seis meses, com calendário, mercado e retornos como essenciais, e nesta configuração os insumos não contribuíram.

\paragraph{RQ2.} No \textit{backtest} cego de vinte decisões mensais, a política obteve 2\,186 reais por saca contra 2\,028 da venda imediata e 2\,001 da venda parcelada, ganhos de 7,8 e 9,2\,\%, significativos a 5\,\%. Em quatro safras, de 2022 a 2026, o ganho foi de 6,1 e 9,0\,\%, também significativo, e a política capturou 94\,\% do valor do oráculo, sendo a melhor estratégia não oracular em três das quatro safras. A ação de \textit{hedge} nunca foi escolhida no \textit{backtest}, porque o mercado esteve em \textit{backwardation}, mas passa a ser recomendada na data mais recente, com a curva real por vencimento, para aversão a risco moderada. A simulação por safra, com uma única decisão na colheita, mostrou o limite da política: uma venda total em 1º de julho de 2024 perdeu a alta inteira daquela safra, e nas duas safras cegas a venda parcelada ficou à frente.

\paragraph{RQ3.} A política é insensível à taxa de desconto entre 8 e 15\,\% ao ano, à tarifa de armazenagem entre 1,50 e 5,00 reais por saca ao mês e a choques de base, porque a deriva prevista nas datas em que o motor armazena é uma ordem de grandeza maior que o custo de carregamento. O único parâmetro que altera a política é a aversão a risco: até $\lambda = 0{,}25$ ela preserva a receita e reduz a frequência de armazenar; a partir de $\lambda = 0{,}5$ converge para a venda imediata.

Em síntese, a hipótese do trabalho foi confirmada nas condições estudadas: uma previsão que acerta a direção em seis meses, mesmo errando a magnitude e o curto prazo, gera valor econômico quando acoplada a uma regra de decisão que respeita custos e incerteza. A contribuição principal não é uma arquitetura mais acurada, mas a demonstração de que a avaliação de um modelo preditivo para decisões de comercialização precisa ir além da métrica de erro, e de que a mesma rede pode ser considerada fraca por uma tabela de RMSE e útil por um \textit{backtest} de receita.

\section{Limitações}

As limitações mais importantes foram tratadas no \autoref{sec:ameacas} e são resumidas aqui. A amostra de avaliação é pequena, vinte decisões cegas e trinta e oito no total, em um período de preços recordes, e os intervalos de confiança do ganho são largos. A arquitetura e a dispersão residual foram selecionadas no conjunto de validação, e a semente do modelo final é favorável, de modo que as métricas de validação são otimistas. O horizonte de 21 dias não tem valor preditivo, e a política depende dele como primeira porta de decisão, o que produziu a venda prematura de julho de 2024. A variável de geada é inerte, e o choque climático mais relevante do período de treino não está representado nos dados de entrada. Os custos de armazenagem e quebra são estimativas sem documento de referência, ainda que a sensibilidade mostre que não alteram a decisão. O \textit{hedge} é modelado sem margem de garantia nem ajuste diário, com um contrato estrangeiro como \textit{proxy} do doméstico e, no \textit{backtest}, com curva plana. E o preço do produtor é aproximado pelo indicador CEPEA/ESALQ, sem a base local de cada cooperativa.

\section{Trabalhos futuros}

Os resultados apontam extensões em três direções.

Na previsão, a mais promissora é prever o excedente sobre a curva de futuros em vez do retorno absoluto, deixando à rede apenas o que o mercado ainda não precifica, e ponderar os horizontes na interpolação pela acurácia direcional de cada um. A variável de geada deve ser redefinida, com limiar de 5\,$^{\circ}$C ou como anomalia de temperatura em relação à climatologia do mês, e variáveis de oferta ainda não coletadas, como os estoques certificados da bolsa e as estimativas de safra da CONAB e do USDA, devem ser incorporadas. Uma avaliação em janela rolante ao longo de toda a série, retreinando a rede a cada safra, responderia à ameaça de regime de mercado.

Na decisão, a venda parcelada deve entrar no motor como ação, e não apenas como referência: vender uma fração e reavaliar o restante preservaria a opção de participar de altas não previstas, que foi o que a venda total na colheita de 2024 destruiu. A compra de opções de venda, que limita a perda sem travar o ganho, é a extensão natural da ação de \textit{hedge}, e exige modelar prêmio e volatilidade implícita. O \textit{hedge} em si deve incorporar margem de garantia e ajuste diário.

Na ferramenta, o passo seguinte é uma aplicação móvel ou web que entregue ao produtor, mensalmente, a ação do dia, o gatilho de venda por mês e a classificação de margem sobre o custo de produção, com os parâmetros econômicos e a aversão a risco configuráveis. A estrutura do motor é agnóstica à cultura, e a extensão a outras commodities armazenáveis com mercado a termo, como o milho e a soja, exige novas fontes, novas variáveis e nova calibração, mas não um novo método.
